% answers_short.tex —— 1977 年·上海卷·理科 · 简短答案速查
% 由 tools/make_short_answers.py 自 answers.tex 的 \gksolhead/\gkanswer 提取，
% 并与 1977/1977_数学真题_上海理_答案要点.md 逐题比对后生成，勿手改（改源 MD 后重跑）。

\begin{gknotice}
\textbf{使用说明}：本节只给最终答案便于快速核对，完整推导见解析版。答案为编写者独立推导，\textbf{非官方标准答案}。
\end{gknotice}

\gkpart{题目}
\begin{itemize}[leftmargin=2.6em]
  \item[\textbf{1.}] (1) $\dfrac{b-a}{a+b}$；(2) $-\dfrac12$；(3) $i$ 是该方程的解；(4) 见解析版
  \item[\textbf{2.}] $DE=\dfrac{ab}{a+b}$。
  \item[\textbf{3.}] $BC=\sqrt6$，$\angle C=45^\circ$。
  \item[\textbf{4.}] (1) 见解析版；(2) $6\sqrt7\ \mathrm{cm^2}$；(3) $45^\circ$
  \item[\textbf{5.}] $[-4,-2)\cup(3,4]$。
  \item[\textbf{6.}] $x^2+4y^2=16$，即 $\dfrac{x^2}{16}+\dfrac{y^2}{4}=1$，为椭圆；严格按斜率定义，$A(-4,0)$、$B(4,0)$ 两点不在实际轨迹中。
  \item[\textbf{7.}] 两腰均为 $15$；下底 $\dfrac{45}{2}$，上底 $\dfrac{15}{2}$。此时面积最大，为 $\dfrac{225\sqrt3}{2}$。
  \item[\textbf{8.}] $k=12$，方程组的相同解为 $(x,y)=(6,38)$。
\end{itemize}
